\documentclass[11pt]{article}

\usepackage[letterpaper,margin=1in]{geometry}
\usepackage[T1]{fontenc}
\usepackage[utf8]{inputenc}
\usepackage{lmodern}
\usepackage{microtype}
\usepackage{amsmath,amssymb,mathtools,bm}
\usepackage{slashed}
\usepackage{enumitem}
\usepackage{xcolor}
\usepackage{booktabs}
\usepackage{array}
\usepackage{hyperref}
\usepackage[most]{tcolorbox}
\usepackage{fancyhdr}
\usepackage{titlesec}
\usepackage{parskip}
\usepackage{needspace}

\definecolor{qftblue}{HTML}{245A8D}
\definecolor{qftgreen}{HTML}{2F6B4F}
\definecolor{qftred}{HTML}{A23B3B}
\definecolor{qftgray}{HTML}{F4F6F8}

\hypersetup{
  colorlinks=true,
  linkcolor=qftblue,
  urlcolor=qftblue,
  citecolor=qftgreen,
  pdftitle={PHYS 662 Lecture 14: One-Particle States and Free Spinor Solutions},
  pdfauthor={Nima Lashkari}
}

\pagestyle{fancy}
\fancyhf{}
\lhead{PHYS 662: Quantum Field Theory I}
\rhead{Lecture 14}
\cfoot{\thepage}
\setlength{\headheight}{14pt}

\titleformat{\section}{\Large\bfseries\color{qftblue}}{\thesection.}{0.6em}{}
\titleformat{\subsection}{\large\bfseries\color{qftgreen}}{\thesubsection}{0.6em}{}
\titleformat{\subsubsection}{\normalsize\bfseries\color{qftred}}{\thesubsubsection}{0.6em}{}

\setlist[itemize]{leftmargin=2em,itemsep=0.25em,topsep=0.4em}
\setlist[enumerate]{leftmargin=2em,itemsep=0.25em,topsep=0.4em}

\newcommand{\dd}{\mathrm{d}}
\newcommand{\R}{\mathbb{R}}
\newcommand{\Lag}{\mathcal{L}}
\newcommand{\cM}{\mathcal{M}}
\newcommand{\SO}{\mathrm{SO}}
\newcommand{\OO}{\mathrm{O}}

\newtcolorbox{keyidea}[1][]{
  enhanced, breakable, colback=qftgray, colframe=qftblue,
  boxrule=0.8pt, arc=1.5mm, title={#1}, fonttitle=\bfseries,
  coltitle=white
}
\newtcolorbox{remarkbox}[1][]{
  enhanced, breakable, colback=white, colframe=qftgreen,
  boxrule=0.7pt, arc=1.5mm, title={#1}, fonttitle=\bfseries,
  coltitle=white
}
\newtcolorbox{warningbox}[1][]{
  enhanced, breakable, colback=white, colframe=qftred,
  boxrule=0.7pt, arc=1.5mm, title={#1}, fonttitle=\bfseries,
  coltitle=white
}


\newcommand{\Th}{\Theta}
\newcommand{\Eul}{\mathcal E}
\newcommand{\Lie}{\mathscr L}

\newcommand{\C}{\mathbb C}
\newcommand{\Z}{\mathbb Z}
\newcommand{\SU}{\mathrm{SU}}
\newcommand{\UU}{\mathrm{U}}
\newcommand{\GL}{\mathrm{GL}}
\newcommand{\SL}{\mathrm{SL}}
\newcommand{\Sp}{\mathrm{Sp}}
\newcommand{\tr}{\operatorname{tr}}
\newcommand{\Ad}{\operatorname{Ad}}
\newcommand{\ad}{\operatorname{ad}}
\newcommand{\lieg}{\mathfrak g}
\newcommand{\lieh}{\mathfrak h}
\newcommand{\su}{\mathfrak{su}}
\newcommand{\so}{\mathfrak{so}}
\newcommand{\slie}{\mathfrak{sl}}


\newcommand{\cH}{\mathcal H}
\newcommand{\Spin}{\operatorname{Spin}}
\newcommand{\id}{\mathrm{id}}
\newcommand{\Lor}{\SO^+(1,3)}
\newcommand{\eps}{\varepsilon}
\newcommand{\one}{\mathbf 1}

\newcommand{\smu}{s}
\newcommand{\tmu}{\widetilde s}
\newcommand{\Eps}{E}
\newcommand{\pslash}{\slashed p}
\newcommand{\ISO}{\mathrm{ISO}}
\newcommand{\pL}{P_L}
\newcommand{\pR}{P_R}
\newcommand{\hc}{\mathrm{h.c.}}
\newcommand{\diag}{\operatorname{diag}}
\newcommand{\vol}{\mathrm{vol}}
\newcommand{\supp}{\operatorname{supp}}
\begin{document}
\pagenumbering{gobble}
\begin{titlepage}
\centering
\vspace*{1.2in}
{\scshape\Large PHYS 662 --- Quantum Field Theory I\par}
\vspace{0.35in}
{\Huge\bfseries\color{qftblue} Lecture 14\par}
\vspace{0.15in}
{\LARGE\bfseries One-Particle States\par}
\vspace{0.06in}
{\LARGE\bfseries and Free Spinor Solutions\par}
\vspace{0.16in}
{\large\color{qftgreen} Little Groups, Spin, and the Dirac Equation\par}
\vspace{0.6in}
{\large Nima Lashkari\par}
{\large Fall 2026\par}
\vfill
\begin{minipage}{0.86\textwidth}\small
The Lorentz matrices acting on field components are finite dimensional
and usually nonunitary. Physical one-particle states, by contrast,
belong to a unitary representation of the Poincar\'e group. Momentum
orbits and their stabilizers explain how these statements fit together.
We derive the massive and massless spin labels, then construct the
Dirac wavefunctions that connect covariant fields to particle states.

\medskip
Course page: \url{https://nimalashkari.com/quantum-field-theory-2026-part-i/}
\end{minipage}
\vspace{0.5in}
\end{titlepage}
\pagenumbering{roman}
\begingroup\small\tableofcontents\endgroup
\newpage
\pagenumbering{arabic}

\section*{Roadmap}
\addcontentsline{toc}{section}{Roadmap}
Lecture 11 classified Lorentz representations, Lecture 12 constructed
their spinor bilinears, and Lecture 13 derived the local free actions. We now ask what the observers in different inertial
frames agree to call the same particle.
\begin{enumerate}
\item Construct one-particle Hilbert spaces from momentum orbits.
\item Derive the little group and the Wigner transformation of spin labels.
\item Classify ordinary massive spin and massless helicity representations.
\item Solve, normalize, and sum the free Dirac plane-wave spinors.
\end{enumerate}
The first part concerns quantum states, independently of a particular
field equation. The second shows how the Dirac equation realizes a
massive spin-\(1/2\) representation. We only preview canonical field
quantization to explain the meaning of the negative-frequency modes.

\begin{keyidea}[Conventions and the distinction of spaces]
We use \(\eta=\diag(1,-1,-1,-1)\),
\(p\cdot x=E t-\bm p\cdot\bm x\), and
\(E_{\bm p}=\sqrt{\bm p^2+m^2}>0\).
The matrices of Lecture 13 are
\[
 \gamma^\mu=\begin{pmatrix}0&\widetilde s^\mu\\s^\mu&0\end{pmatrix},
 \quad s^\mu=(\one,\bm\sigma),\quad
 \widetilde s^\mu=(\one,-\bm\sigma),\quad
 \bar u=u^\dagger\gamma^0.
\]
Write \(U\) for a unitary operator on states and \(\mathcal S\) for
a finite-dimensional matrix on spinor components. They act on
different spaces. We use \(U(a)=e^{-ia_\mu P^\mu}\) for translations
of states; reversing the convention for \(a\) changes only this phase.
\end{keyidea}

\section{What is the representation-theoretic definition of a particle?}
\subsection{States form a representation space}
In a relativistic quantum theory, spacetime symmetries act on the
Hilbert space by unitary operators for the connected Poincar\'e group,
or, more precisely, its covering group when half-integer spin is
allowed. A stable elementary one-particle sector is modeled by a
positive-energy irreducible unitary representation of that group.
If there are internal multiplets or several species, the total
one-particle space may contain several such representations.

\emph{Irreducible} means that the representation space has no nontrivial
closed invariant subspace. A particular state vector is not itself
``an irreducible representation.'' A localized wavepacket, a momentum
eigenstate, and a spin superposition can all be states of the same
particle species. Momentum eigenstates are generalized states; their
delta-function normalization does not make them normalizable
Hilbert-space vectors.

The translation generators commute:
\begin{equation}
 [P^\mu,P^\nu]=0,\qquad
 P^\mu|p,a\rangle=p^\mu|p,a\rangle,
 \qquad U(b)|p,a\rangle=e^{-ib\cdot p}|p,a\rangle.
 \label{eq:translation}
\end{equation}
Here \(a\) is a residual polarization label. Lorentz transformations
move the eigenvalue from \(p\) to \(\Lambda p\). For example, with
our convention,
\begin{equation}
 U(\Lambda)^{-1}P^\mu U(\Lambda)
    =\Lambda^\mu{}_{\nu}P^\nu
 \quad\Longrightarrow\quad
 P^\mu U(\Lambda)|p,a\rangle
    =(\Lambda p)^\mu U(\Lambda)|p,a\rangle.
 \label{eq:momentumcov}
\end{equation}
The polarization label can also change. Determining that change is
the role of the little group.

\subsection{A nonrelativistic comparison}
For a free nonrelativistic particle,
\(H=\bm P^2/(2m)\). A momentum eigenstate is a character of the
translation group, but it is not an invariant sector of the full
Galilei group: boosts shift \(\bm p\) by a quantity proportional to
\(m\bm v\). At rest, all rotations preserve the momentum, so spin
is a representation of the rotation group, with cover \(\SU(2)\).

At a nonzero momentum, rotations about \(\bm p\) are only the
rotation-only stabilizer. They are not the full group controlling
massive spin: a general rotation can be accompanied by a boost that
returns the momentum to its original value. The relativistic little
group makes this distinction precise. Inspecting only rotations
about a generic momentum would incorrectly discard most of a massive
spin multiplet.

\section{Momentum orbits and the invariant inner product}
\subsection{Mass shells}
Lorentz transformations preserve \(p^2\) and, on timelike or null
orbits, the sign of the nonzero energy. For a massive positive-energy
particle, the orbit is
\begin{equation}
 \mathcal O_m^+=\{p:p^2=m^2,\ p^0>0\},\qquad m>0.
 \label{eq:massshell}
\end{equation}
Every point on it is a boost of \(k=(m,0,0,0)\). For a massless
particle, the positive light cone with its tip removed is the orbit;
we may choose \(k=(\kappa,0,0,\kappa)\), where \(\kappa>0\) is
an arbitrary reference energy, not a mass or an additional particle label.

The operator \(P_\mu P^\mu\) commutes with the Poincar\'e action and
is constant on an irreducible sector. Its value is \(m^2\).
Our physical classification restricts the spectrum to the future
timelike or null cone. Spacelike orbits cannot maintain an energy
bounded below under all boosts. The zero-momentum orbit, which
contains the invariant vacuum representation, is separate from the
nonzero massless particle orbit.

\subsection{Deriving the measure}
Start with a manifestly invariant measure on the positive mass shell:
\begin{align}
 \dd^4p\,\delta(p^2-m^2)\theta(p^0)
 &=\dd^3\bm p\,\dd p^0\,
   \delta\big((p^0)^2-E_{\bm p}^2\big)\theta(p^0)\nonumber\\
 &=\frac{\dd^3\bm p}{2E_{\bm p}}.
 \label{eq:measurederive}
\end{align}
The factor \(2E_{\bm p}\) is the derivative of the delta-function
argument at its positive root. Recall that for a function $f(x)$ with simple poles at $x_i$ we have
\begin{align}
    \delta(f(x))=\sum_i \frac{\delta(x-x_i)}{|f'(x_i)|}
\end{align}
In our case, we have
\begin{align}
    \delta(p^2-m^2)\Theta(p^0)=\frac{\delta(p^0-E_p)+\delta(p^0+E_p)}{2E_p}\Theta(p^0)=\frac{\delta(p^0-E_p)}{2E_p}\Theta(p^0)
\end{align}



We include Fourier factors by defining
\begin{equation}
 \dd\Pi_p=\frac{\dd^3\bm p}{(2\pi)^3\,2E_{\bm p}}.
 \label{eq:measure}
\end{equation}
The factor \((2\pi)^3\) is conventional; the factor \(E_{\bm p}\)
is required by Lorentz invariance.

Choose covariant normalization,
\begin{equation}
 \langle p,a|q,b\rangle
   =(2\pi)^3\,2E_{\bm p}\,\delta^3(\bm p-\bm q)\delta_{ab}.
 \label{eq:ketnorm}
\end{equation}
A one-particle wavepacket can then be expanded as
\begin{equation}
 |f\rangle=\sum_a\int\dd\Pi_p\,f_a(p)|p,a\rangle,
 \qquad
 \langle f|f\rangle=\sum_a\int\dd\Pi_p\,|f_a(p)|^2.
 \label{eq:wavepacket}
\end{equation}
For a scalar particle there is one component. For spin, there is a
fiber of polarization states at each momentum. Locally this looks
like an \(L^2\) space of vector-valued functions; for massless helicity,
choices of polarization basis over the momentum sphere may require
more than one patch. The invariant content does not depend on that
choice of basis.

\section{The little group and the Wigner transformation}
\subsection{Bring a momentum back to the reference momentum}
Fix a reference point \(k\) on the orbit. For each \(p\), choose a
standard Lorentz transformation \(L(p)\) such that
\begin{equation}
 L(p)k=p,\qquad |p,a\rangle=U(L(p))|k,a\rangle.
 \label{eq:standardboost}
\end{equation}
For a massive particle we can choose rotation-free boosts. The
\emph{little group} is the stabilizer
\begin{equation}
 G_k=\{W:Wk=k\}.
 \label{eq:littledef}
\end{equation}
Apply a general \(\Lambda\) to a state produced from \(k\). The
combination
\begin{equation}
 \boxed{W(\Lambda,p)=L(\Lambda p)^{-1}\Lambda L(p)}
 \label{eq:wigner}
\end{equation}
returns \(k\) to itself, so it belongs to \(G_k\). This formula is
the missing bridge between a Lorentz transformation and its action
on a finite set of spin labels.

If the reference polarization states carry a unitary irreducible
representation \(D\) of \(G_k\), then
\begin{equation}
 \boxed{U(\Lambda)|p,a\rangle
   =\sum_b|\Lambda p,b\rangle
        D_{ba}\!\left(W(\Lambda,p)\right).}
 \label{eq:ketwigner}
\end{equation}
There is no square-root energy factor with normalization
\eqref{eq:ketnorm}. With states normalized instead to
\(\delta^3(\bm p-\bm q)\), a compensating energy factor would
appear. The two conventions should not be mixed.

\subsection{Why the induced action is unitary}
Changing variables to \(p'=\Lambda p\) in the wavepacket integral
gives
\begin{equation}
 (U(\Lambda)f)_b(p')
    =\sum_a D_{ba}\!\left(W(\Lambda,\Lambda^{-1}p')\right)
                  f_a(\Lambda^{-1}p').
 \label{eq:wffransform}
\end{equation}
The measure is invariant and \(D^\dagger D=\one\), so
\(\|U(\Lambda)f\|^2=\|f\|^2\). The full representation is
infinite dimensional because the momentum runs over an entire orbit,
even when \(D\) has only finitely many polarization components.

Composition is consistent because
\begin{equation}
 W(\Lambda_2\Lambda_1,p)
  =W(\Lambda_2,\Lambda_1p)W(\Lambda_1,p).
 \label{eq:cocycle}
\end{equation}
This follows by inserting
\(L(\Lambda_1p)L(\Lambda_1p)^{-1}\) into \eqref{eq:wigner}.
Changing the standard transformations by
\(L(p)\mapsto L(p)R(p)\), with \(R(p)\in G_k\), merely changes
the momentum-dependent polarization basis. It conjugates the Wigner
matrices at the initial and final momenta without changing the
particle representation.

\begin{keyidea}[Field spin and particle spin answer different questions]
The Lorentz labels \((j_L,j_R)\) classify finite-dimensional field
components. The little-group labels classify physical polarizations
at fixed momentum. Equations of motion and, for gauge fields, gauge
equivalence connect the two. A four-component Dirac field does not
describe four independent particle spin polarizations of one charge.
\end{keyidea}

\section{Massive particles: the full rotation group}
\subsection{The rest-frame stabilizer}
For \(k=(m,0,0,0)\), any transformation satisfying \(\Lambda k=k\)
leaves the time direction fixed. The Lorentz condition then forces
the spatial block to be an orthogonal three-by-three matrix; in the
connected proper group it has determinant \(+1\). Thus
\begin{equation}
 G_k=\SO(3),\qquad \widetilde G_k=\SU(2).
 \label{eq:massivelittle}
\end{equation}
Its irreducible unitary representations are labeled by
\begin{equation}
 j=0,\tfrac12,1,\tfrac32,\ldots,\qquad a=-j,-j+1,\ldots,j.
 \label{eq:massivespin}
\end{equation}
There are \(2j+1\) spin states for each momentum and for each
independent particle species. A Dirac particle has two, and its
antiparticle has another two. A massive Majorana particle has two
with no independent antiparticle species.

The little group at a moving momentum is conjugate to the rest-frame
group, \(G_p=L(p)G_kL(p)^{-1}\). Its elements generally combine
boosts and rotations. It is therefore still a three-dimensional
rotation group, not just the \(\SO(2)\) of rotations about \(\bm p\).
Successive non-collinear boosts lead to a nontrivial Wigner rotation
of the spin basis; this is why spin labels need not remain unchanged
under a boost.

\subsection{The second invariant}
Besides \(P^2\), the Poincar\'e algebra has the Pauli--Lubanski
invariant. With orientation chosen so that \(\epsilon^{0123}=1\),
define
\begin{equation}
 \mathcal W^\mu=-\frac12\epsilon^{\mu\nu\rho\sigma}
                        P_\nu M_{\rho\sigma}.
 \label{eq:PL}
\end{equation}
Our overall sign makes \(\mathcal W^\mu=(0,m\bm J)\) in the rest
frame when \(J_i=\tfrac12\eps_{ijk}M_{jk}\). It follows that
\begin{equation}
 P^2=m^2,\qquad
 \mathcal W^2=-m^2j(j+1),\qquad P_\mu\mathcal W^\mu=0.
 \label{eq:casimirs}
\end{equation}
The sign of \(\mathcal W\) itself is conventional; its square and
the massive spin label are not. The two invariants summarize the
massive spacetime classification, apart from additional internal
quantum numbers.

\section{Massless particles: helicity and the null little group}
\subsection{The stabilizer is larger than rotations about momentum}
Choose \(k=(\kappa,0,0,\kappa)\). Using the rotation--boost algebra
of Lecture 11 and its boost sign convention, the infinitesimal
generators preserving this vector are
\begin{equation}
 J=J_3,\qquad N_1=K_1+J_2,\qquad N_2=K_2-J_1.
 \label{eq:nullgens}
\end{equation}
For example, a boost generated by \(K_1\) and a rotation generated
by \(J_2\) change the transverse component of \(k\) in opposite
ways. Their sum leaves \(k\) fixed. Direct use of the commutators gives
\begin{equation}
 [J,N_1]=iN_2,\qquad [J,N_2]=-iN_1,\qquad [N_1,N_2]=0.
 \label{eq:iso2}
\end{equation}
This is the Lie algebra of the Euclidean group of the plane,
\(\ISO(2)=\SO(2)\ltimes\R^2\), whose translations are the two
null rotations. They are not ordinary spacetime translations;
they are elements of the Lorentz stabilizer of a null momentum.
For spinorial states one uses the corresponding little group inside
the Lorentz covering group.

\subsection{Why ordinary finite-polarization particles have helicity}
In a finite-dimensional unitary little-group representation,
\(N_1,N_2\) are commuting Hermitian matrices, so their joint
eigenvalues form a finite set in \(\R^2\). Conjugation by
\(e^{-i\theta J}\) rotates that set. A nonzero vector has a whole
circle as its rotation orbit, which cannot fit into a finite set.
Therefore
\begin{equation}
 N_1=N_2=0
 \label{eq:trivialnull}
\end{equation}
for every finite-dimensional unitary representation of this type.
The remaining irreducible rotation representation is one dimensional:
\begin{equation}
 D(\theta)=e^{-ih\theta},\qquad
 h=\bm J\cdot\widehat{\bm p}.
 \label{eq:helicity}
\end{equation}
The number \(h\) is the helicity. Rotations have period \(4\pi\)
in the Lorentz double cover, so \(h\in\tfrac12\Z\). Requiring an
ordinary representation of the Lorentz group itself would restrict
\(h\) to integers. For helicity representations
\(\mathcal W^\mu=hP^\mu\) with our sign choice, and
\(\mathcal W^2=0\). Thus \(P^2\) and \(\mathcal W^2\) alone
do not distinguish the different helicities at zero mass.

If nontrivial null-rotation representations are allowed, unitary
representations with infinitely many polarizations also exist:
the continuous-spin representations. Lorentz symmetry alone does
not rule them out. The usual photon and massless fermion discussion
restricts attention to the finite-polarization helicity sectors.

\subsection{Helicity pairs and the role of discrete symmetries}
For the connected Poincar\'e group, a fixed helicity is already an
irreducible representation. A photon with both \(h=+1\) and
\(h=-1\) polarizations contains two such sectors. Parity reverses
momentum while preserving the axial vector \(\bm J\), and therefore
exchanges \(h\leftrightarrow-h\). A parity-invariant theory must
include the appropriate partner; connected Lorentz symmetry alone
does not require it. Particle--antiparticle assignments and CPT give
further relations once a local quantum field theory is specified.

Massive particles also admit helicity eigenstates, but helicity is
not invariant under all boosts: an observer can overtake a massive
particle and reverse its momentum. For a massless helicity
representation, proper orthochronous transformations preserve the
helicity label. Chirality is instead a property of the field's
Lorentz representation. For massless on-shell Weyl solutions these
notions are related; for massive Dirac solutions they must not be
identified.

In our explicit matrices a positive-frequency upper Weyl solution
obeys \((E-\bm\sigma\cdot\bm p)u=0\) and has \(h=+1/2\).
A lower Weyl solution has \(h=-1/2\). A convention that exchanges
the two chiral names exchanges this verbal assignment. The actual
helicity eigenvalue follows from the equation and the rotation
generator \(\bm\sigma/2\).

\section{Free Dirac solutions at fixed momentum}
\subsection{Positive- and negative-frequency equations}
Consider the free equation
\begin{equation}
 (i\gamma^\mu\partial_\mu-m)\Psi=0,\qquad m>0.
 \label{eq:dirac}
\end{equation}
Since its coefficients are constant, Fourier modes diagonalize its
spacetime dependence. Use the same future-directed on-shell momentum
\(p=(E_{\bm p},\bm p)\) in both ans\"atze:
\begin{align}
 \Psi^{(+)}(x)&=u_s(p)e^{-ip\cdot x},
       & (\slashed p-m)u_s(p)&=0,\nonumber\\
 \Psi^{(-)}(x)&=v_s(p)e^{+ip\cdot x},
       & (\slashed p+m)v_s(p)&=0,
 \label{eq:uvdef}
\end{align}
where \(\slashed p=\gamma^\mu p_\mu\). The columns \(u_s,v_s\) are ordinary
complex coefficient spinors, not Grassmann fields. In the quantum field
expansion, the mode operators will carry the fermionic statistics.
The negative-frequency mode
has Fourier four-momentum \(-p\); \(p\) remains the positive-energy
label used for its eventual antiparticle interpretation.
The general solution is a superposition of these modes, not a single
plane wave. Squaring the equation gives the mass shell, but one still
has to solve the spinor constraints in \eqref{eq:uvdef}.

\subsection{Rest-frame solutions and their spin labels}
At \(k=(m,\bm0)\), the two equations reduce to
\(\gamma^0u=u\) and \(\gamma^0v=-v\). Let
\(\xi_s,\eta_s\), \(s=1,2\), each be an orthonormal basis of
two-component columns. Then
\begin{equation}
 u_s(k)=\sqrt m\begin{pmatrix}\xi_s\\\xi_s\end{pmatrix},\qquad
 v_s(k)=\sqrt m\begin{pmatrix}\eta_s\\-\eta_s\end{pmatrix}.
 \label{eq:restspinors}
\end{equation}
The \(\sqrt m\) gives \(u^\dagger u=v^\dagger v=2m\) at rest.
For example \(\xi_1=(1,0)^T\), \(\xi_2=(0,1)^T\) are spin-up
and spin-down along \(z\). Spatial rotations act with the same
\(\sigma_i/2\) on both blocks, displaying the two-dimensional
massive little-group representation directly.

The \(\eta_s\) basis can be chosen independently. If one wants
\(v_s=C\bar u_s^T\), the compatible choice with Lecture 13 is
\(\eta_s=E\xi_s^*\). This convention may relabel the spin of the
two-component column relative to a basis with \(\eta_s=\xi_s\).
Completeness sums are independent of this basis choice.

\subsection{Boosting the rest spinors}
Set
\begin{equation}
 A_p=\sqrt{E_{\bm p}\one+\bm p\cdot\bm\sigma},\qquad
 B_p=\sqrt{E_{\bm p}\one-\bm p\cdot\bm\sigma},
 \label{eq:AB}
\end{equation}
where the square roots are the positive Hermitian ones. For \(m>0\)
the eigenvalues \(E_{\bm p}\pm|\bm p|\) are positive, so these
matrices are well defined. They commute and obey
\begin{equation}
 A_p^2=E+\bm p\cdot\bm\sigma,\quad
 B_p^2=E-\bm p\cdot\bm\sigma,\quad A_pB_p=m\one.
 \label{eq:ABids}
\end{equation}
An explicit form, useful without diagonalization, is
\begin{equation}
 A_p=\frac{(E+m)\one+\bm p\cdot\bm\sigma}{\sqrt{2(E+m)}},\qquad
 B_p=\frac{(E+m)\one-\bm p\cdot\bm\sigma}{\sqrt{2(E+m)}}.
 \label{eq:ABexplicit}
\end{equation}
Squaring the numerator verifies these formulas using
\((\bm p\cdot\bm\sigma)^2=\bm p^2\one\).

Write \(E=m\cosh\eta\), \(\bm p=m\sinh\eta\,\bm n\), with
\(|\bm n|=1\). Then
\(A_p=\sqrt m\,e^{+\eta\bm n\cdot\bm\sigma/2}\) and
\(B_p=\sqrt m\,e^{-\eta\bm n\cdot\bm\sigma/2}\).
The standard boost taking rest to \(+\bm p\) has rapidity
\(\bm\rho=-\eta\bm n\) in our active-boost convention. Thus the
spinor boost is \(\diag(A_p/\sqrt m,B_p/\sqrt m)\), and
\begin{equation}
 \boxed{u_s(p)=\begin{pmatrix}A_p\xi_s\\B_p\xi_s\end{pmatrix},
 \qquad v_s(p)=\begin{pmatrix}A_p\eta_s\\-B_p\eta_s\end{pmatrix}.}
 \label{eq:uv}
\end{equation}
This also verifies the equations without appealing to the boost:
\begin{equation}
 \slashed p=\begin{pmatrix}0&A_p^2\\B_p^2&0\end{pmatrix},
 \qquad A_p^2B_p=mA_p,\quad B_p^2A_p=mB_p.
 \label{eq:slashblocks}
\end{equation}

\section{Normalization, orthogonality, and spin sums}
\subsection{Dirac inner products and ordinary inner products}
Because \(A_p,B_p\) are Hermitian,
\begin{equation}
 \bar u_s=(\xi_s^\dagger B_p,\xi_s^\dagger A_p),\qquad
 \bar v_s=(-\eta_s^\dagger B_p,\eta_s^\dagger A_p).
 \label{eq:adjointsuv}
\end{equation}
Their Lorentz-scalar inner products are
\begin{equation}
 \bar u_su_t=2m\delta_{st},\qquad
 \bar v_sv_t=-2m\delta_{st},\qquad
 \bar u_sv_t=\bar v_su_t=0.
 \label{eq:uvbarnorm}
\end{equation}
For example \(\bar u_su_t=\xi_s^\dagger(B_pA_p+A_pB_p)\xi_t\).
Their ordinary positive norms are instead
\begin{equation}
 u_s^\dagger u_t=2E\delta_{st},\qquad
 v_s^\dagger v_t=2E\delta_{st}.
 \label{eq:uvdagnorm}
\end{equation}
Both follow from \(A_p^2+B_p^2=2E\one\).

It is important to specify the momentum when claiming orthogonality:
\begin{equation}
 u_s^\dagger(p)v_t(p)
    =2\xi_s^\dagger(\bm p\cdot\bm\sigma)\eta_t
 \label{eq:notorth}
\end{equation}
need not vanish. The negative-energy eigenvector of the
\emph{same spatial-momentum Hamiltonian} \(h_D(\bm p)\) is
\(v_t(-\bm p)\), not \(v_t(\bm p)\). Since
\(A_{-\bm p}=B_{\bm p}\),
\begin{equation}
 u_s^\dagger(\bm p)v_t(-\bm p)=0.
 \label{eq:trueorth}
\end{equation}
This agrees with Hermitian Hamiltonian eigenvectors of distinct
eigenvalues being orthogonal.

\subsection{Completeness sums}
Use \(\sum_s\xi_s\xi_s^\dagger=\sum_s\eta_s\eta_s^\dagger=\one_2\).
Block multiplication then gives
\begin{align}
 \sum_su_s(p)\bar u_s(p)
  &=\begin{pmatrix}m\one&A_p^2\\B_p^2&m\one\end{pmatrix}
    =\boxed{\slashed p+m},\label{eq:usum}\\
 \sum_sv_s(p)\bar v_s(p)
  &=\begin{pmatrix}-m\one&A_p^2\\B_p^2&-m\one\end{pmatrix}
    =\boxed{\slashed p-m}.\label{eq:vsum}
\end{align}
These identities replace a sum over spin labels by a covariant
matrix, which is why they are central to later scattering calculations.
For a fixed future-directed on-shell \(p\), the projectors onto the
\(\slashed p\) eigenspaces are
\begin{equation}
 \Pi_+(p)=\frac{\slashed p+m}{2m},\qquad
 \Pi_-(p)=\frac{m-\slashed p}{2m}.
 \label{eq:slashprojectors}
\end{equation}
They obey \(\Pi_\pm^2=\Pi_\pm\), \(\Pi_+\Pi_-=0\), and
\(\Pi_++\Pi_-=1\). In particular \(\sum v\bar v/(2m)\) is
\(-\Pi_-\), not \(\Pi_-\), because \(\bar vv=-2m\).
These projectors need not be Hermitian in the ordinary Euclidean
component inner product.

For a fixed spatial-momentum Hamiltonian the corresponding orthogonal
projectors are instead
\begin{equation}
 Q_\pm(\bm p)=\frac12\left(1\pm\frac{h_D(\bm p)}{E}\right),
 \quad Q_+=\frac{1}{2E}\sum_su_s(\bm p)u_s^\dagger(\bm p),
 \quad Q_-=\frac{1}{2E}\sum_sv_s(-\bm p)v_s^\dagger(-\bm p).
 \label{eq:hamprojectors}
\end{equation}
The distinction resolves the apparent conflict between
\eqref{eq:notorth} and Hamiltonian orthogonality.

\section{How the wavefunctions connect fields to particles}
\subsection{A preview of the quantum field expansion}
With the normalization used above, the free Dirac field is expanded as
\begin{equation}
 \Psi(x)=\sum_s\int\dd\Pi_p\,
 \left[b_s(p)u_s(p)e^{-ip\cdot x}
          +d_s^\dagger(p)v_s(p)e^{+ip\cdot x}\right].
 \label{eq:fieldexpansion}
\end{equation}
The operators \(b_s^\dagger(p)\) and \(d_s^\dagger(p)\) create
particle and antiparticle states, respectively. For example,
\(\{b_s(p),b_t^\dagger(q)\}=(2\pi)^3 2E_{\bm p}
\delta_{st}\delta^3(\bm p-\bm q)\), with the same normalization
for \(d\). The remaining independent anticommutators vanish.
Thus the coefficient of a negative-frequency mode is an antiparticle
\emph{creation} operator, rather than a negative-energy particle
annihilation operator.

After the vacuum contribution is removed by normal ordering, the
free Hamiltonian takes the form
\begin{equation}
 :H:=\sum_s\int\dd\Pi_p\,E_{\bm p}
          \left[b_s^\dagger(p)b_s(p)+d_s^\dagger(p)d_s(p)\right].
 \label{eq:positiveH}
\end{equation}
The spectrum of excitations is positive. Deriving this formula and
the equal-time canonical algebra in detail belongs to field
quantization; its role here is to explain why both sets of solutions
in \eqref{eq:uvdef} are needed. For a Majorana field the reality
condition relates the two mode sectors, so there is no independent
antiparticle creation operator.

\subsection{The intertwiner between the two representations}
The boosted columns \(u_s(p)\) are not just a list of solutions.
They connect the little-group spin index \(s\) with the covariant
Dirac component index. With the same standard boosts used to define
the states,
\begin{equation}
 \mathcal S(\Lambda)u_s(p)
   =\sum_tu_t(\Lambda p)
      D^{(1/2)}_{ts}\!\left(W(\Lambda,p)\right).
 \label{eq:intertwiner}
\end{equation}
To derive it, write
\(u_s(p)=\mathcal S(L(p))u_s(k)\), insert
\(\Lambda L(p)=L(\Lambda p)W(\Lambda,p)\), and use the
ordinary spin-\(1/2\) action of the little group on the rest columns.
This explains how a nonunitary matrix on four field components can
coexist with a unitary transformation of the two particle spin states.
The momentum dependence of the wavefunctions supplies the bridge.

\section*{Summary}
\addcontentsline{toc}{section}{Summary}
A one-particle representation is assembled from a Lorentz momentum
orbit and a unitary little-group representation. The invariant
mass-shell measure and Wigner matrices make the full action unitary.
For massive particles the little group is \(\SU(2)\), giving
\(2j+1\) polarizations. For ordinary finite-polarization massless
particles its null-rotation part acts trivially, leaving helicity.
The boosted Dirac wavefunctions realize the massive spin-\(1/2\)
case; their normalizations and spin sums follow from simple
two-by-two matrix identities. Quantization turns their two frequency
sectors into positive-energy particles and antiparticles.

\newpage
\section*{Questions for review}
\addcontentsline{toc}{section}{Questions for review}
\begin{enumerate}
\item Why is a momentum eigenstate not by itself a Poincar\'e
irreducible representation?
\item Derive the mass-shell measure, including the factor \(2E\).
\item Prove that \(W(\Lambda,p)\) fixes \(k\) and obeys
\eqref{eq:cocycle}.
\item Prove unitarity of \eqref{eq:wffransform} using the measure and
the little-group inner product.
\item Explain why the moving massive little group is not merely
rotations about the momentum axis.
\item Verify the algebra \eqref{eq:iso2}. Which assumption forces
\(N_1=N_2=0\), and what happens when it is removed?
\item Why can a theory contain both photon helicities while a fixed
helicity is already irreducible under the connected group?
\item Check the square-root matrices \eqref{eq:ABexplicit} and use
them to verify both Dirac equations in \eqref{eq:uvdef}.
\item Compare \(\bar u v\), \(u^\dagger(p)v(p)\), and
\(u^\dagger(p)v(-\bm p)\). Why do they have different answers?
\item Derive the two spin sums and explain the sign of the
negative-eigenvalue projector.
\item Explain which step replaces negative-frequency classical
solutions with positive-energy antiparticle states.
\item Derive \eqref{eq:intertwiner} from the rest-frame spinors and
the definition of the Wigner transformation.
\end{enumerate}

\section*{Sources and reading}
\addcontentsline{toc}{section}{Sources and reading}
The main source is the supplied \emph{Lecture 13: Single Particle
States}, pp.\ 8--21. Its introductory spinor review is incorporated
in course Lectures 12 and 13. The standard-boost construction and the
finite-polarization qualification make the particle-classification
argument explicit. For a systematic account of momentum orbits and
little groups, see X.~Bekaert and N.~Boulanger,
\emph{The unitary representations of the Poincar\'e group in any
spacetime dimension},
\href{https://arxiv.org/abs/hep-th/0611263}{arXiv:hep-th/0611263}.
For spinor conventions, compare H.~K.~Dreiner, H.~E.~Haber, and
S.~P.~Martin, \href{https://arxiv.org/abs/0812.1594}{arXiv:0812.1594}.
\end{document}
